% Generated draft: requires mathematical and source review.
\documentclass[11pt]{article}
\usepackage[margin=25mm]{geometry}
\usepackage{fontspec}
\setmainfont{Latin Modern Roman}
\usepackage{amsmath,amssymb,mathtools,unicode-math}
\setmathfont{Latin Modern Math}
\usepackage{enumitem,needspace,xurl,hyperref}
\setlength{\parindent}{0pt}
\setlength{\parskip}{5pt}
\setlength{\emergencystretch}{3em}
\newcommand{\N}{\mathbb{N}}
\newcommand{\Z}{\mathbb{Z}}
\newcommand{\Q}{\mathbb{Q}}
\newcommand{\R}{\mathbb{R}}
\newcommand{\C}{\mathbb{C}}
\newcommand{\F}{\mathbb{F}}
\newcommand{\A}{\mathbb{A}}
\newcommand{\PP}{\mathbb{P}}
\newcommand{\E}{\mathbb{E}}
\newcommand{\eps}{\varepsilon}
\newcommand{\dd}{\,\mathrm{d}}
\begin{document}
\section*{UnsolvedMath ID: 1200030}
\textbf{Source problem code:} catalogue entry
\subsection{Definitions and mathematical statement}
Let $T_3=(V,E)$ be the infinite $3$-regular tree. Under the usual full-automorphism convention, a factor of i.i.d. process with values in a finite alphabet $A$ is a random labeling $Y\in A^V$ obtained as $Y=F(U)$, where $(U_v)_{v\in V}$ are independent uniform $[0,1]$ labels and $F:[0,1]^V\to A^V$ is measurable and $\operatorname{Aut}(T_3)$-equivariant, modulo null sets. Write $\mathcal F_A$ for the set of its laws. Weak convergence of laws on $A^V$ means convergence of all finite-coordinate distributions; equivalently, for every finite $W\subset V$ and $z\in A^W$,
\[
 \nu_m\{y:y|_W=z\}\longrightarrow\nu\{y:y|_W=z\}.
\]

\textbf{Closure question.} Is $\mathcal F_A$ closed for this topology? More explicitly, must $\nu\in\mathcal F_A$ whenever $\nu_m\in\mathcal F_A$ and $\nu_m\Rightarrow\nu$? The source does not specify the output alphabet or acting group; the fixed-finite-alphabet and full-automorphism conventions above make one standard interpretation explicit.

\textbf{Majority-on-balls question.} Let $B_n(v)=\{w\in V:d(v,w)\leq n\}$. In the usual binary interpretation, start with independent unbiased labels $\xi_w\in\{-1,1\}$ and form a process whose value at $v$ is the majority sign of
\[
 \sum_{w\in B_n(v)}\xi_w.
\]
After specifying an equivariant rule for ties, denote its law by $\nu_n$. The source asks whether \emph{the weak limit} of these majority processes is a factor of i.i.d., that is, whether the intended limit $\nu_n\Rightarrow\nu$ belongs to $\mathcal F_{\{-1,1\}}$. The record neither specifies the input distribution and tie rule nor establishes existence or uniqueness of this limit, and it does not specify a subsequence. These are unresolved qualifications rather than additional conclusions. In particular, $|B_n(v)|=3\cdot2^n-2$ is even for $n\geq1$, so ties require a convention at finite radii.

The supplied summary labels both questions OPEN-TRIAGE, connects them with closures of factor-of-i.i.d. laws and with rounding or belief-propagation limits, and reports no settled answer within a capped search. Its classification is expressly provisional and is not an independently verified openness claim.
\subsection{Short English statement}
Are weak limits of factor-of-i.i.d. processes on the $3$-regular tree still factors of i.i.d.? The specific test case is the putative limit of majority labels on larger balls. The source leaves the alphabet, symmetry convention, tie rule, and precise limit unspecified.
\subsection{Sources}
\begin{itemize}
\item Local UnsolvedMath registry, permanent ID 1200030, source code catalogue entry, ``Some Questions --- Question 30''; status open, classification OPEN-TRIAGE, updated 2026-08-06.
\item Abert, \emph{Some questions}, list dated 2010, Question 30, \url{https://www.renyi.hu/~abert/questions.pdf}. Reported access: 2026-07-29; extraction: manual transcription. The background says the 2010 author list was inspected but current status was not independently verified. Status and rights: NEEDS\_REVIEW; public author-hosted PDF.
\item The supplied research summary names Ab\'ert--Cs\'oka--Herrero--Lipnowitz--Vervloet, Benjamini--Schramm, and Bowen in connection with this program, without particular publications or theorem references. Its search cap and provisional status assessment are retained.
\end{itemize}
\end{document}
